\naslov{Valovna enačba na realni osi}

Imamo podan koeficient $c > 0$.
Gledamo enačbo
\[
  u_{tt} - c^2 u_{xx} = 0,
\]
kjer je $x \in (a,b)$ in $t > 0$.
To je hiperbolična enačba, ki ima kanonično obliko $w_{\xi \eta} = 0$, pri čemer
sta novi koordinati oblike $\xi = x - ct$ in $\eta = x + ct$.
Če integriramo $w_{\xi \eta} = 0$, dobimo $w = D(\xi) + E(\eta)$ oziroma $u(x,
t) = F(x - ct) + G(x + ct)$ za poljubni $F, G \in \zvodvi{2}$.
Predpis deluje tudi, če funkciji nista odvedljivi.
Če sta le odsekoma zvezni, obstajata zaporedji gladkih funkcij $F_n$ in $G_n$,
ki konvergirata k $F$ in $G$ po točkah in enakomerno na kompaktih, kjer sta $F$
in $G$ zvezni.
Potem bo zaporedje $u_n = F_n(x - ct) + G(x+ct)$ na enak način konvergiralo k
$u$.

\begin{definicija}
  Če sta $F, G$ odsekoma zvezni na $\R$, se $u(t,x) = F(x-ct) + G(x+ct)$ imenuje
  \pojem{posplošena rešitev} za valovno enačbo.
\end{definicija}

\vprasanje{Kakšne oblike so rešitve valovne enačbe? Kaj pa posplošene rešitve?}